\section{Causal Probabilistic Transformers}
\label{sec:model}

% Mirrors Wu & Tu Section 2 ("Probabilistic Transformers"): the model first, then inference,
% then variants. Their Section 2 runs about 2.3 pages of a two-column ACL page; budget the
% same here. Figure 1 (the factor graph) belongs in 2.2, early, exactly as theirs does.

We first recall the model of \citet{wu2023probabilistic} only as far as we need it
(\Cref{sec:recap}), then state the obstruction that prevents a direct causal version
(\Cref{sec:obstruction}), then give the model (\Cref{sec:chain}), its output mechanism
(\Cref{sec:output}), its inference procedure (\Cref{sec:inference}), and finally some variants.

\subsection{Background}
\label{sec:recap}

\FROMPDF{Wu \& Tu \S2.1--\S2.2 --- variable set, the two potential functions, MFVI}
\TODO{compress to about half a column. Only what \Cref{sec:obstruction,sec:output} actually
use: the label variables $\Zv{i}$ over $\nlab$ labels; the head variables $\Hv{i}{c}$ over
candidate heads, one set per channel; the unary factor $\phi_u(\Zv{i}) = \exp(\Sfac_{w_i,\Zv{i}})$
with $\Sfac \in \R^{|\Vocab| \times \nlab}$; the ternary factor over $(\Hv{i}{c}, \Zv{i}, \Zv{j})$
with score matrix $\Tfac{c}$; and that MFVI unrolled for $\niter$ iterations has the shape of a
$\niter$-layer transformer encoder with shared weights. Note explicitly that the word $w_i$
enters \emph{only} through $\Sfac$ --- \Cref{sec:output} turns on that fact.}

\subsection{Why a causal mask is not enough}
\label{sec:obstruction}

\TODO{about half a column. Wu \& Tu's paper contains no theorem environments; keep this to one
theorem statement in prose-adjacent form and push the proof to \Cref{app:proof}.}

The obvious route to a causal model is to mask the ternary factors so that $\Hv{i}{c}$ may only
point to positions $j \le i$. This does not give a causal model. Informally, three properties
that one would want cannot hold together:

\begin{itemize}
  \item \textbf{Causality:} the posterior at position $i$ depends only on $w_{1:i}$;
  \item \textbf{A single global energy:} one energy function over the sentence, normalised once;
  \item \textbf{Contextuality:} the representation at $i$ genuinely depends on its prefix.
\end{itemize}

\begin{theorem}[Trilemma]
\label{thm:trilemma}
\TODO{formal statement} No model satisfies all three simultaneously.
\end{theorem}

The mechanism is the partition function: a factor coupling positions $i<j$ contributes to the
marginal at $i$ whenever it contributes to the marginal at $j$, and forbidding the first while
keeping the second is inconsistent with a single normalisation. Masking the factors changes
which configurations have non-zero score; it does not stop the surviving ones from being
normalised jointly. The proof is in \Cref{app:proof}.

Causality and contextuality are what the model is for, so the single global energy is what we
give up.

\begin{remark}[Never formed, not deleted]
\label{rmk:never-formed}
Building the symmetric model and then removing the anti-causal contribution gives a different
object from building a model in which that contribution never exists. This distinction is what
\Cref{sec:chain} implements and what \Cref{sec:exp-correctness} tests.
\end{remark}

\subsection{A directed chain of conditional CRFs}
\label{sec:chain}

\FROMPDF{Part~II of the main document --- the factorisation and the conditioning statement}

\begin{figure}[t]
\centering
\TODO{Figure 1. Two panels of the \emph{same} factor graph: (a) encoder, $\Wv{i}$ shaded
(observed); (b) decoder, $\Wv{i}$ unshaded (latent), with the prefix boxed as a frozen
condition. Same factors, same edges, different shading. This single figure carries the claim
that encoding and decoding are two conditioning patterns of one model, and it is the analogue
of Wu \& Tu's Figure 1, which is also introduced at exactly this point.}
\caption{The factor graph of the causal probabilistic transformer for $\seqlen = 3$.}
\label{fig:factor-graph}
\end{figure}

Instead of one CRF over the sentence, we define a chain of conditional CRFs, one per position.
The CRF at position $t$ ranges over $(\Zv{t}, \Hv{t}{1..\nchan}, \Wv{t})$ and is conditioned on
the prefix posteriors $\{\QZ{s}\}_{s<t}$, which enter as fixed numbers:
%
\begin{equation}
\label{eq:chain}
  \TODO{the factorisation}
\end{equation}

\paragraph{The prefix is a frozen condition.}
Each $\QZ{s}$, $s<t$, is data at step $t$: no message travels from $t$ back to $s$. The
anti-causal term of \Cref{thm:trilemma} is therefore never formed, and causality holds by
construction rather than by masking.

\paragraph{Frozen forward, not frozen backward.}
\emph{Frozen} means constant with respect to step $t$'s inference problem, not detached in
autodiff. Training gradients flow backwards through the whole prefix computation, exactly as
they do through cached activations in a causal transformer; this is how $\Sfac$ is trained in
both of its roles (\Cref{sec:output}). \NOTE{say this explicitly in the paper. It is the single
easiest thing for a reader to get backwards, and getting it backwards makes the model look
either non-causal or untrainable.}

\subsection{The output mechanism}
\label{sec:output}

This is the part of the construction we most want checked, and we present it as such.

\paragraph{The design we rejected.}
The obvious way to obtain next-token logits is to project $\QZ{t}$ through a fresh matrix in
$\R^{\nlab \times |\Vocab|}$. Such a matrix corresponds to no factor in the graph: the output
would not traverse the factor graph at all, and the model would be a CRF with a neural head
attached. We report this because rejecting it is what forced the mechanism below.

\paragraph{Unclamping the word variable.}
The word variable is already in the graph and $\Sfac$ is already a binary factor; in the encoder
it is evaluated with its word argument clamped to the observed token. The output mechanism is to
stop clamping it:
%
\begin{equation}
\label{eq:unclamp}
  \TODO{$\readout{t}(a) \propto \dots$ in the notation of \texttt{causal\_pt\_output\_note.pdf}}
\end{equation}
%
No factor is added and no parameter is created. Next-token prediction is the posterior over a
variable the model already had.

\paragraph{Consequences.}
Three properties follow, and we stress that they are consequences rather than design choices:

\begin{enumerate}
  \item \textbf{Tying is forced.} One factor $\Sfac$ mediates both directions, so input and
        output word embeddings are the same matrix. Untying is not a hyperparameter; it is a
        different model.
  \item \textbf{The output bias is the word unary.} \TODO{state}
  \item \textbf{One model, two conditioning patterns.} The encoder clamps $\Wv{i}$; the decoder
        unclamps it and directs the chain. Same parameters, same factors
        (\Cref{fig:factor-graph}).
\end{enumerate}

\subsection{Inference}
\label{sec:inference}

\NOTE{Resolved 2026-08-05: \S17.2 proposes MFVI as the mainline but \S23.3 supersedes it. The
exact readout is the model; the mean-field readout is the ablation of \Cref{sec:exp3}. Do not
present them as two equal options.}

\FROMPDF{Part~III \S16--\S17.1 --- update equations and schedule}
\begin{align}
\label{eq:update-Z}
  \QZ{t} &\propto \TODO{} \\
\label{eq:update-H}
  \QH{t}{c} &\propto \TODO{}
\end{align}

\paragraph{Exact readout.}
The subgraph at position $t$ is a star centred on $\Zv{t}$, hence a tree, so sum-product is
exact: no approximation and no iteration are involved.

\begin{proposition}
\label{prop:exact-readout}
\FROMPDF{\S23.3 --- confirm the form and the seeding by $\rfac{c}$}
\begin{equation}
\label{eq:exact-readout}
  \log \readout{t}(a) = \sum_{c=1}^{\nchan} \LSE_{j \in \Dep{t}} \Bfac{c}_{j,a},
\end{equation}
seeded with the root key $\rfac{c}_a$.
\end{proposition}

Along the sequence $\Dep{t}$ grows by one position per step, so \Cref{eq:exact-readout} is a
causal prefix log-sum-exp: one \texttt{logcumsumexp} scan, $O(\seqlen\nlab\nchan)$, fully
parallel over positions. The exact readout therefore costs neither the layer-parallel schedule
nor an inner loop --- with it there is no query stream and no predictive iteration at all. Two
costs are real and we state them rather than bury them: the $|\Vocab| \times \nlab$ readout is a
log-sum-exp rather than a matmul (same FLOPs, worse hardware constants, chunked over the
vocabulary with a fused cross-entropy), and LSE-pooling gradients are untested at
language-model scale.

\paragraph{Every step is valid inference.}
\FROMPDF{Part~III \S18 Check~3, as restated in \S23.3}
Mean field where exactness is impossible (the chain), exact sum-product where the subgraph is a
tree (the slot).

\subsection{Extensions and variants}
\label{sec:variants}

% Wu & Tu keep this subsection short and push the rest to their Appendix B. Do the same.

\paragraph{Distance.} \TODO{the causal half of the clipped ternary potential: only $i-j>0$
exists, so there are $\clip+1$ buckets rather than $2\clip+1$.}

\paragraph{Update schedule.} \TODO{parallel (a depth-$\niter$ shared causal transformer) vs.\
serial ($\tau$ inner rounds per slot); say which produced the reported numbers.}

\paragraph{Global variables.} \TODO{the Appendix B.3 globals $\Gv{t}$, proposed by
\citet{wu2023probabilistic} as an in-graph substitute for the missing feed-forward structure and
never tested there. We treat them as a measured variable (\Cref{sec:exp-ablations}), not an
assumption.}

Further variants are in \Cref{app:variants}.
